% Wspólna część obu pism: preambuła, dane nadawcy, uzasadnienie i załączniki.
% Pisma: pismo-burmistrz.tex i pismo-zkm.tex. Kompilacja (dwa razy):
%   xelatex pismo-burmistrz.tex
%   xelatex pismo-zkm.tex
% Załącznik nr 1 (zalacznik-1-zmiany.tex) tworzy narzędzie zkmplan.py:
%   cp ../data/plan/<data>/zmiany.tex zalacznik-1-zmiany.tex
% Załączniki nr 2 i 3 (raporty do druku) tworzy print_reports.py:
%   python3 print_reports.py --pdf   (pliki w ../reports/druk/)
\documentclass[11pt,a4paper]{article}
\usepackage[a4paper, top=2.2cm, bottom=2.4cm, left=2.5cm, right=2.2cm]{geometry}
\usepackage{fontspec}
\setmainfont{texgyretermes}[Extension=.otf, UprightFont=*-regular, BoldFont=*-bold,
  ItalicFont=*-italic, BoldItalicFont=*-bolditalic]
\usepackage{babel}
\babelprovide[import, main]{polish}
\usepackage{microtype}
\usepackage{array, booktabs, longtable}
\usepackage{enumitem}
\usepackage{parskip}
\usepackage{pdfpages}
\usepackage[hidelinks]{hyperref}
\emergencystretch=3em
\setlist{itemsep=3pt, topsep=4pt, leftmargin=1.6em}
\pagestyle{plain}

% ---- Dane do uzupełnienia przed wydrukiem --------------------------------
\newcommand{\Nadawca}{Marcin Jasiukowicz}
\newcommand{\AdresNadawcy}{\makebox[5.5cm]{\dotfill}\\ 14-200 Iława}
\newcommand{\EmailNadawcy}{contact@yasiu.pl}
\newcommand{\Data}{\makebox[2.8cm]{\dotfill}}
\newcommand{\RaportPrzesiadki}{https://claude.ai/artifact/SjGg6GoZq5NishuoizYrvL}
\newcommand{\RaportKorekty}{https://claude.ai/artifact/LHpHerFfAXQ8om2pxWeYN2}
\newcommand{\DrukPrzesiadki}{../reports/druk/zalacznik-2-przesiadki.pdf}
\newcommand{\DrukKorekty}{../reports/druk/zalacznik-3-korekty.pdf}
% ---------------------------------------------------------------------------

\newcommand{\Czesc}[1]{\par\medskip{\large\bfseries #1}\par\smallskip}
\newcommand{\Tytul}[1]{\par\smallskip{\bfseries #1}\par}

% Nagłówek: nadawca po lewej, miejscowość i data po prawej, adresat niżej po prawej.
\newcommand{\Naglowek}[1]{%
\thispagestyle{empty}
\begin{minipage}[t]{0.52\textwidth}
\Nadawca\\
\AdresNadawcy\\
e-mail: \EmailNadawcy
\end{minipage}\hfill
\begin{minipage}[t]{0.42\textwidth}\raggedleft
Iława, dnia \Data{} 2026 r.
\end{minipage}

\vspace{1.1cm}
\hfill\begin{minipage}[t]{0.55\textwidth}
#1
\end{minipage}
\vspace{0.9cm}}

\newcommand{\Wstep}{%
Przeanalizowałem rozkład jazdy autobusów ZKM Iława obowiązujący od 1 września 2026 r. na tle rozkładu jazdy pociągów na stacji Iława Główna (dane PKP Polskich Linii Kolejowych S.A.). Wyniki i metoda są opisane w raportach załączonych do pisma (załączniki nr~2 i~3).}

\newcommand{\Uzasadnienie}{%
\Tytul{Uzasadnienie}

Autobusy ZKM nie są obecnie dopasowane do pociągów. W dzień roboczy autobus odjeżdża z przystanku Dworcowa w chwili, gdy podróżni idą jeszcze z peronu, przy 50\,\% przyjazdów pociągów. Rozkład autobusów ułożony losowo dałby 47\,\%. Liczba dogodnych przesiadek również odpowiada rozkładowi losowemu.

Dogodna przesiadka to czekanie najwyżej 10 minut po dojściu z peronu na przystanek (4 minuty) albo z przystanku na peron (5 minut). Do przyjazdu pociągu doliczono typowe opóźnienie jego kategorii: 3 minuty dla pociągów PKP Intercity. Liczone są pary pociąg--linia, ponieważ podróżny potrzebuje konkretnej linii.

Proponowane zmiany wyznaczyło narzędzie, które przesuwa kursy o najwyżej 5 minut, a kursy kursujące tylko w dni nauki szkolnej o najwyżej 2 minuty. Przesunięty kurs zmienia godziny na wszystkich przystankach o tyle samo. Narzędzie zachowuje obiegi autobusów, przerwy kierowców wymagane ustawą o czasie pracy kierowców, liczbę autobusów i łączny czas pracy kierowców. Kursy każdej linii zachowują kolejność. Obiegi i służby narzędzie odtworzyło z rozkładu, ponieważ nie są one publikowane. Każdy etap zachowuje zmiany etapu poprzedniego, więc rozkład można zmieniać stopniowo.

\begin{center}
\begin{tabular}{@{}lrrrrl@{}}
\toprule
 & \multicolumn{4}{c}{Dogodne przesiadki dziennie} & \\
\cmidrule(lr){2-5}
Typ dnia & obecnie & etap 1 & etap 2 & wszystkie zmiany & liczba zmian \\
\midrule
Dzień roboczy (188 kursów) & 148 & 180 (+22\,\%) & 199 (+35\,\%) & 244 (+65\,\%) & 20 / 40 / 87 \\
Sobota (89 kursów) & 91 & 117 (+29\,\%) & 132 (+45\,\%) & 133 (+46\,\%) & 20 / 40 / 43 \\
Niedziela (54 kursy) & 46 & 70 (+52\,\%) & -- & 76 (+64\,\%) & 20 / 28 \\
\bottomrule
\end{tabular}
\end{center}

Najwięcej zyskują osoby dojeżdżające autobusem na pociąg oraz podróżni pociągów regionalnych. Przy rzeczywistych przyjazdach pociągów PKP Intercity z 13 dni roboczych liczba dogodnych przesiadek z pociągu do autobusu rośnie o 9\,\%.

Na stacji Iława Główna wsiada i wysiada średnio 4700 osób na dobę (Urząd Transportu Kolejowego, 2025). Jeżeli autobusem ZKM jedzie dalej od 5\,\% do 15\,\% z nich, krótsze czekanie na przesiadkę daje w dzień roboczy od 1,3\,\% do 4,5\,\% więcej przejazdów autobusem do dworca i z dworca. W ciągu roku to od około 950 do 10\,000 przejazdów.

Narzędzie znalazło w dzień roboczy 8 okresów, w których autobus stoi na przystanku Dworcowa dłużej, niż trwa kurs tam i z powrotem. Dodatkowe kursy w tym czasie dałyby łącznie około 25 dogodnych przesiadek dziennie. Przy obiegach ze zmianą linii na przystanku Dworcowa obecny rozkład wymaga w szczycie 15 autobusów zamiast 18.

Rozkład w formacie GTFS pozwoli pokazać połączenia autobus--pociąg w wyszukiwarkach połączeń i łatwo sprawdzać skutki każdej zmiany rozkładu. Rozkład ZKM w tym formacie przygotowałem na potrzeby analizy. Jestem gotów przekazać narzędzie, dane i rozkład w formacie GTFS oraz spotkać się z przedstawicielami Urzędu Miasta i ZKM.}

\newcommand{\Podpis}{%
\vspace{0.8cm}
\hfill\begin{minipage}{0.42\textwidth}\centering
Z poważaniem\\[1.4cm]
\makebox[5.5cm]{\dotfill}\\
\Nadawca
\end{minipage}}

\newcommand{\MaterialyIZalaczniki}{%
\Tytul{Materiały w internecie}
\begin{itemize}
\item Raport „Przesiadki Iława Główna”: \url{\RaportPrzesiadki}
\item Raport „Korekty rozkładu ZKM Iława”: \url{\RaportKorekty}
\end{itemize}

\Tytul{Załączniki}
\begin{enumerate}
\item Wykaz proponowanych zmian rozkładu jazdy (dzień roboczy, sobota, niedziela).
\item Raport „Przesiadki Iława Główna” (analiza rozkładów i opóźnień), wersja do druku.
\item Raport „Korekty rozkładu ZKM Iława” (propozycja korekt), wersja do druku.
\end{enumerate}}

% Załącznik nr 1 i raporty (załączniki nr 2 i 3), jeżeli ich pliki PDF istnieją.
\newcommand{\Zalaczniki}{%
\newpage
\section*{Załącznik nr 1. Wykaz proponowanych zmian rozkładu jazdy}

Wykaz dotyczy obiegów, w których autobus zostaje na jednej linii. Godziny są godzinami na przystanku Dworcowa: obecna $\to$ \textbf{proponowana}.
\begin{itemize}[itemsep=1pt]
\item \textbf{Etap}: 1 -- najskuteczniejsze zmiany (do 20 kursów), 2 -- do 40 kursów łącznie, 3 -- pozostałe zmiany. Każdy etap zachowuje zmiany poprzedniego.
\item \textbf{Inne przyst.}: zmiana godzin na pozostałych przystankach kursu, w minutach.
\item \textbf{Zysk}: liczba dogodnych przesiadek dziennie, którą daje ta jedna zmiana.
\item \textbf{Pociągi}: pociągi, z których (\emph{z}) lub na które (\emph{na}) przesiadka staje się dogodna. Kursywą: pociągi, dla których przesiadka przestaje być dogodna. Zmiana towarzysząca zachowuje przerwę kierowcy, postój autobusu lub kolejność kursów po innej zmianie.
\end{itemize}

{\footnotesize
\setlength{\tabcolsep}{3pt}
\renewcommand{\arraystretch}{1.15}
% Generated by zkmplan.py. Do not edit: run the tool again.
% Generated 2026-09-29T15:47:24, ZKM feed 4aad5903cd
\subsection*{Dzień roboczy}
Plan na dzień 30.09.2026, sprawdzony na 20 dniach tego typu (30.09.2026--26.10.2026). Dobre przesiadki dziennie (pary pociąg--linia, z zapasem na typowe opóźnienie pociągu): obecnie 148.

\begin{tabular}{@{}lrrr@{}}\toprule
Etap & Zmienione kursy (razem) & Dobre przesiadki dziennie & Zmiana \\ \midrule
1 & 20 & 180 & +22\,\% \\
2 & 40 & 199 & +35\,\% \\
3 & 87 & 244 & +65\,\% \\
\bottomrule\end{tabular}

\begin{longtable}{@{}c c >{\raggedright}p{2.4cm} l l r r >{\raggedright\arraybackslash}p{4.3cm}@{}}\toprule
Etap & Linia & Kierunek & Przyjazd & Odjazd & Inne & Zysk & Pociągi \\
 & & & \multicolumn{2}{c}{na Dworcowej} & przyst. & /dzień & (dobra przesiadka po zmianie) \\ \midrule\endhead
\bottomrule\endfoot
1 & 3 & Księżnej Dobrawy & -- & 9:03\,$\to$\,\textbf{9:00} & $-$3 & +1.9 & z R 90652 (08:46); z R 95623 (08:46) \\
1 & 1 & Dworcowa & 11:17\,$\to$\,\textbf{11:12} & -- & $-$5 & +1.8 & na TLK 8328 (11:17); na IC 6115 (11:17) \\
1 & 2 & Iława ul. Lipowy Dwór (ogrody) & -- & 13:00\,$\to$\,\textbf{13:05} & +5 & +1.9 & z IC 1815 (12:58) \\
1 & 3 & Dworcowa & 14:45\,$\to$\,\textbf{14:48} & -- & +3 & +2.0 & na EIP 3553 (15:03); na R 55205 (15:03) \\
1 & 1 & Dworcowa & 14:46\,$\to$\,\textbf{14:48} & -- & +2 & +2.0 & na EIP 3553 (15:03); na R 55205 (15:03) \\
1 & 1 & Długa & -- & 15:02\,$\to$\,\textbf{15:05} & +3 & +2.0 & z IC 1515 (14:57); z IC 8114 (14:58) \\
1 & 3 & Nowa Wieś & -- & 15:16\,$\to$\,\textbf{15:14} & $-$2 & +2.0 & z IC 1515 (14:57); z IC 8114 (14:58) \\
1 & 3 & Dworcowa & 15:44\,$\to$\,\textbf{15:48} & -- & +4 & +2.0 & na R 90647 (16:01); na EIP 3505 (16:03) \\
1 & 1 & Dworcowa & 15:47\,$\to$\,\textbf{15:48} & -- & +1 & +1.0 & na EIP 3505 (16:03) \\
1 & 2 & Dworcowa & 16:20\,$\to$\,\textbf{16:24} & -- & +4 & +1.8 & na IC 1614 (16:36); na R 90225 (16:39) \\
1 & 1 & Dworcowa & 16:22\,$\to$\,\textbf{16:25} & -- & +3 & +1.0 & na R 90225 (16:39) \\
1 & 1 & Dworcowa & 19:00\,$\to$\,\textbf{18:58} & -- & $-$2 & +2.0 & na IC 5110 (19:03); na EIP 4503 (19:03) \\
1 & 3 & Dworcowa & 19:00\,$\to$\,\textbf{18:58} & -- & $-$2 & +2.0 & na IC 5110 (19:03); na EIP 4503 (19:03) \\
1 & 2 & Iława ul. Lipowy Dwór (ogrody) & -- & 19:08\,$\to$\,\textbf{19:10} & +2 & +2.0 & z IC 5110 (19:01); z EIP 4503 (19:02) \\
1 & 3 & Dworcowa & 19:13\,$\to$\,\textbf{19:12} & -- & $-$1 & +0.9 & na IC 6119 (19:17) \\
1 & 2 & Iława ul. Lipowy Dwór (ogrody) & -- & 19:23\,$\to$\,\textbf{19:24} & +1 & +0.9 & z IC 6119 (19:16) \\
1 & 3 & Księżnej Dobrawy & -- & 20:42\,$\to$\,\textbf{20:41} & $-$1 & +0.9 & z R 90246 (20:27) \\
1 & 2 & Dworcowa & 20:49\,$\to$\,\textbf{20:48} & -- & $-$1 & +1.0 & na EIC 5100 (20:53) \\
1 & 1 & Długa & -- & 21:04\,$\to$\,\textbf{21:07} & +3 & +2.0 & z TLK 35105 (20:58); z R 90113 (21:03) \\
1 & 3 & Nowa Wieś & -- & 21:38\,$\to$\,\textbf{21:37} & $-$1 & +0.9 & z IC 7509 (21:20) \\
2 & 2 & Dworcowa & 5:53\,$\to$\,\textbf{5:52} & -- & $-$1 & +0.9 & na R 90250 (05:57) \\
2 & 2 & Iława ul. Lipowy Dwór (ogrody) & -- & 6:05\,$\to$\,\textbf{6:04} & $-$1 & +0.9 & z R 90650 (05:50) \\
2 & 1 & Długa & -- & 7:24\,$\to$\,\textbf{7:26} & +2 & +1.0 & z R 90240 (07:22) \\
2 & 3 & Dworcowa & 7:47\,$\to$\,\textbf{7:48} & -- & +1 & +1.0 & na EIP 1501 (08:03) \\
2 & 3 & Dworcowa & 8:25\,$\to$\,\textbf{8:26} & -- & +1 & +1.0 & na IC 1618 (08:40) \\
2 & 1 & Dworcowa & 9:12\,$\to$\,\textbf{9:10} & -- & $-$2 & +1.0 & na R 90751 (09:15) \\
2 & 7 & Nowa Wieś & -- & 9:15\,$\to$\,\textbf{9:14} & $-$1 & +1.0 & z IC 1511 (08:57) \\
2 & 3 & Dworcowa & 10:36\,$\to$\,\textbf{10:38} & -- & +2 & +1.0 & na EIP 5402 (10:53) \\
2 & 3 & Księżnej Dobrawy & -- & 10:36\,$\to$\,\textbf{10:37} & +1 & +1.0 & z IC 5710 (10:29) \\
2 & 4 & Piaskowa & -- & 11:20\,$\to$\,\textbf{11:19} & $-$1 & +1.0 & z EIP 4509 (11:02) \\
2 & 2 & Iława ul. Lipowy Dwór (ogrody) & -- & 11:30\,$\to$\,\textbf{11:29} & $-$1 & +0.9 & z TLK 8328 (11:12) \\
2 & 3 & Nowa Wieś & -- & 11:50\,$\to$\,\textbf{11:52} & +2 & +1.0 & z R 90654 (11:48) \\
2 & 1 & Dworcowa & 12:37\,$\to$\,\textbf{12:38} & -- & +1 & +1.0 & na EIP 5352 (12:53) \\
2 & 2 & Iława ul. Lipowy Dwór (ogrody) & -- & 13:35\,$\to$\,\textbf{13:33} & $-$2 & +0.9 & z IC 7515 (13:16) \\
2 & 2 & Dworcowa & 13:49\,$\to$\,\textbf{13:48} & -- & $-$1 & +1.0 & na EIP 5404 (13:53) \\
2 & 1 & Dworcowa & 14:30\,$\to$\,\textbf{14:29} & -- & $-$1 & +1.0 & na IC 5714 (14:34) \\
2 & 2 & Dworcowa & 16:37\,$\to$\,\textbf{16:38} & -- & +1 & +1.0 & na EIP 5308 (16:53) \\
2 & 3 & Nowa Wieś & -- & 17:15\,$\to$\,\textbf{17:14} & $-$1 & +1.0 & z TLK 3829 (16:57) \\
2 & 2 & Dworcowa & 17:39\,$\to$\,\textbf{17:38} & -- & $-$1 & +1.0 & na R 55201 (17:43) \\
2 & 2 & Iława ul. Lipowy Dwór (ogrody) & -- & 21:06\,$\to$\,\textbf{21:07} & +1 & +0.9 & z R 90113 (21:03) \\
3 & 1 & Długa & -- & 5:05\,$\to$\,\textbf{5:07} & +2 & +1.0 & z IC 5118 (04:59) \\
3 & 3 & Dworcowa & 5:27\,$\to$\,\textbf{5:31} & -- & +4 & +1.0 & na R 90205 (05:46) \\
3 & 3 & Nowa Wieś & -- & 5:52\,$\to$\,\textbf{5:54} & +2 & +0.9 & z R 90650 (05:50) \\
3 & 3 & Dworcowa & 6:38\,$\to$\,\textbf{6:39} & -- & +1 & +0.0 & \textit{zmiana towarzysząca} \\
3 & 1 & Dworcowa & 6:40\,$\to$\,\textbf{6:45} & -- & +5 & +1.0 & na R 90101 (07:00) \\
3 & 1 & Długa & -- & 6:42\,$\to$\,\textbf{6:47} & +5 & +1.0 & z IC 5708 (06:39) \\
3 & 3 & Nowa Wieś & -- & 6:52\,$\to$\,\textbf{6:54} & +2 & +1.0 & z R 90623 (06:50) \\
3 & 2 & Iława ul. Lipowy Dwór (ogrody) & -- & 7:00\,$\to$\,\textbf{6:56} & $-$4 & +1.0 & z IC 5708 (06:39); z R 90101 (06:44); \textit{traci: z EIP 5300 (06:52)} \\
3 & 2 & Dworcowa & 7:08\,$\to$\,\textbf{7:12} & -- & +4 & +1.0 & na R 90240 (07:27) \\
3 & 2 & Dworcowa & 7:46\,$\to$\,\textbf{7:48} & -- & +2 & +1.0 & na EIP 1501 (08:03) \\
3 & 1 & Dworcowa & 7:49\,$\to$\,\textbf{7:51} & -- & +2 & +0.0 & \textit{zmiana towarzysząca} \\
3 & 1 & Długa & -- & 8:00\,$\to$\,\textbf{8:02} & +2 & +1.0 & z IC 1521 (07:55) \\
3 & 2 & Iława ul. Lipowy Dwór (ogrody) & -- & 8:08\,$\to$\,\textbf{8:10} & +2 & +1.0 & z EIP 1501 (08:02) \\
3 & 3 & Dworcowa & 8:39\,$\to$\,\textbf{8:43} & -- & +4 & +1.0 & na R 90221 (08:57) \\
3 & 2 & Dworcowa & 10:14\,$\to$\,\textbf{10:16} & -- & +2 & +1.0 & na IC 5710 (10:30) \\
3 & 2 & Dworcowa & 10:36\,$\to$\,\textbf{10:38} & -- & +2 & +1.0 & na EIP 5402 (10:53) \\
3 & 4 & Dworcowa & 10:44\,$\to$\,\textbf{10:48} & -- & +4 & +1.0 & na EIP 4509 (11:03) \\
3 & 3 & Księżnej Dobrawy & -- & 10:55\,$\to$\,\textbf{11:00} & +5 & +1.0 & z EIP 5402 (10:52) \\
3 & 3 & Dworcowa & 11:42\,$\to$\,\textbf{11:40} & -- & $-$2 & +0.9 & na R 90242 (11:45) \\
3 & 2 & Dworcowa & 12:00\,$\to$\,\textbf{11:58} & -- & $-$2 & +1.0 & na EIP 3509 (12:03) \\
3 & 4 & Dworcowa & 12:19\,$\to$\,\textbf{12:21} & -- & +2 & +1.0 & na R 90223 (12:36) \\
3 & 2 & Dworcowa & 12:21\,$\to$\,\textbf{12:26} & -- & +5 & +1.0 & na IC 1706 (12:40) \\
3 & 3 & Dworcowa & 12:40\,$\to$\,\textbf{12:45} & -- & +5 & +1.0 & na IC 5114 (12:59) \\
3 & 2 & Dworcowa & 12:56\,$\to$\,\textbf{12:54} & -- & $-$2 & +1.0 & na IC 5114 (12:59) \\
3 & 3 & Dworcowa & 13:46\,$\to$\,\textbf{13:48} & -- & +2 & +1.0 & na EIP 3507 (14:03) \\
3 & 1 & Dworcowa & 13:57\,$\to$\,\textbf{13:55} & -- & $-$2 & +1.0 & na R 90200 (14:00) \\
3 & 3 & Dworcowa & 14:31\,$\to$\,\textbf{14:29} & -- & $-$2 & +1.0 & na IC 5714 (14:34) \\
3 & 4 & Dworcowa & 14:58\,$\to$\,\textbf{14:54} & -- & $-$4 & +1.0 & na IC 8114 (14:59) \\
3 & 2 & Dworcowa & 15:17\,$\to$\,\textbf{15:12} & -- & $-$5 & +1.0 & na IC 7107 (15:17) \\
3 & 3 & Księżnej Dobrawy & -- & 15:20\,$\to$\,\textbf{15:24} & +4 & +1.0 & z IC 7107 (15:16) \\
3 & 3 & Księżnej Dobrawy & -- & 15:40\,$\to$\,\textbf{15:42} & +2 & +1.0 & z R 90244 (15:38) \\
3 & 2 & Dworcowa & 15:56\,$\to$\,\textbf{15:53} & -- & $-$3 & +0.9 & na R 90209 (15:58) \\
3 & 3 & Dworcowa & 16:22\,$\to$\,\textbf{16:24} & -- & +2 & +0.9 & na R 90225 (16:39) \\
3 & 2 & Iława ul. Lipowy Dwór (ogrody) & -- & 16:25\,$\to$\,\textbf{16:30} & +5 & +1.0 & z R 90202 (16:26) \\
3 & 1 & Długa & -- & 16:35\,$\to$\,\textbf{16:38} & +3 & +1.0 & z R 90625 (16:34) \\
3 & 3 & Dworcowa & 16:36\,$\to$\,\textbf{16:38} & -- & +2 & +1.0 & na EIP 5308 (16:53) \\
3 & 3 & Dworcowa & 17:04\,$\to$\,\textbf{17:02} & -- & $-$2 & +1.0 & na TLK 3829 (17:07) \\
3 & 3 & Dworcowa & 17:35\,$\to$\,\textbf{17:38} & -- & +3 & +1.0 & na EIP 5406 (17:53) \\
3 & 2 & Iława ul. Lipowy Dwór (ogrody) & -- & 17:46\,$\to$\,\textbf{17:42} & $-$4 & +1.0 & z R 55201 (17:28) \\
3 & 2 & Dworcowa & 19:53\,$\to$\,\textbf{19:48} & -- & $-$5 & +1.1 & na EIP 5102 (19:53) \\
3 & 1 & Długa & -- & 19:55\,$\to$\,\textbf{19:52} & $-$3 & +0.9 & z R 90660 (19:38) \\
3 & 3 & Księżnej Dobrawy & -- & 19:57\,$\to$\,\textbf{20:00} & +3 & +1.0 & z EIP 5102 (19:52) \\
3 & 2 & Iława ul. Lipowy Dwór (ogrody) & -- & 20:06\,$\to$\,\textbf{20:02} & $-$4 & +1.0 & z R 90227 (19:48) \\
3 & 3 & Dworcowa & 20:43\,$\to$\,\textbf{20:48} & -- & +5 & +1.0 & na EIP 3503 (21:03) \\
3 & 2 & Iława ul. Lipowy Dwór (ogrody) & -- & 20:46\,$\to$\,\textbf{20:41} & $-$5 & +0.9 & z R 90246 (20:27) \\
3 & 2 & Dworcowa & 21:25\,$\to$\,\textbf{21:20} & -- & $-$5 & +1.0 & na R 95636 (21:25) \\
3 & 1 & Dworcowa & 22:42\,$\to$\,\textbf{22:45} & -- & +3 & +1.0 & na IC 1519 (23:00) \\
\end{longtable}

\subsection*{Sobota}
Plan na dzień 3.10.2026, sprawdzony na 4 dniach tego typu (3.10.2026--24.10.2026). Dobre przesiadki dziennie (pary pociąg--linia, z zapasem na typowe opóźnienie pociągu): obecnie 91.

\begin{tabular}{@{}lrrr@{}}\toprule
Etap & Zmienione kursy (razem) & Dobre przesiadki dziennie & Zmiana \\ \midrule
1 & 20 & 117 & +29\,\% \\
2 & 40 & 132 & +45\,\% \\
3 & 43 & 133 & +46\,\% \\
\bottomrule\end{tabular}

\begin{longtable}{@{}c c >{\raggedright}p{2.4cm} l l r r >{\raggedright\arraybackslash}p{4.3cm}@{}}\toprule
Etap & Linia & Kierunek & Przyjazd & Odjazd & Inne & Zysk & Pociągi \\
 & & & \multicolumn{2}{c}{na Dworcowej} & przyst. & /dzień & (dobra przesiadka po zmianie) \\ \midrule\endhead
\bottomrule\endfoot
1 & 3 & Dworcowa & 6:49\,$\to$\,\textbf{6:48} & -- & $-$1 & +1.0 & na EIP 5300 (06:53) \\
1 & 3 & Dworcowa & 7:52\,$\to$\,\textbf{7:51} & -- & $-$1 & +1.0 & na IC 1521 (07:56) \\
1 & 3 & Dworcowa & 8:47\,$\to$\,\textbf{8:48} & -- & +1 & +1.0 & na EIP 1503 (09:03) \\
1 & 1 & Długa & -- & 9:05\,$\to$\,\textbf{9:06} & +1 & +1.0 & z IC 5116 (08:58) \\
1 & 2 & Iława ul. Lipowy Dwór (ogrody) & -- & 9:25\,$\to$\,\textbf{9:24} & $-$1 & +1.0 & z R 90730 (09:10) \\
1 & 2 & Dworcowa & 11:01\,$\to$\,\textbf{11:03} & -- & +2 & +2.0 & na TLK 8328 (11:17); na IC 6115 (11:17) \\
1 & 1 & Dworcowa & 11:37\,$\to$\,\textbf{11:38} & -- & +1 & +1.0 & na EIP 5302 (11:53) \\
1 & 2 & Dworcowa & 12:44\,$\to$\,\textbf{12:48} & -- & +4 & +2.0 & na R 90632 (13:01); na EIP 4507 (13:03) \\
1 & 2 & Iława ul. Lipowy Dwór (ogrody) & -- & 13:01\,$\to$\,\textbf{13:05} & +4 & +2.0 & z IC 1815 (12:58); z IC 5114 (12:58) \\
1 & 1 & Długa & -- & 13:20\,$\to$\,\textbf{13:15} & $-$5 & +3.0 & z IC 1815 (12:58); z IC 5114 (12:58); z EIP 4507 (13:02) \\
1 & 3 & Dworcowa & 14:30\,$\to$\,\textbf{14:29} & -- & $-$1 & +1.0 & na IC 5714 (14:34) \\
1 & 1 & Dworcowa & 14:52\,$\to$\,\textbf{14:53} & -- & +1 & +1.0 & na IC 1515 (15:08) \\
1 & 2 & Dworcowa & 14:52\,$\to$\,\textbf{14:53} & -- & +1 & +1.0 & na IC 1515 (15:08) \\
1 & 1 & Długa & -- & 15:00\,$\to$\,\textbf{15:05} & +5 & +2.0 & z IC 1515 (14:57); z IC 8114 (14:58) \\
1 & 2 & Iława ul. Lipowy Dwór (ogrody) & -- & 16:37\,$\to$\,\textbf{16:38} & +1 & +1.0 & z R 90625 (16:34) \\
1 & 2 & Iława ul. Lipowy Dwór (ogrody) & -- & 18:10\,$\to$\,\textbf{18:09} & $-$1 & +1.0 & z EIP 5406 (17:52) \\
1 & 3 & Nowa Wieś & -- & 19:15\,$\to$\,\textbf{19:14} & $-$1 & +1.0 & z IC 1517 (18:57) \\
1 & 3 & Dworcowa & 20:02\,$\to$\,\textbf{20:00} & -- & $-$2 & +1.0 & na IC 5320 (20:05) \\
1 & 2 & Dworcowa & 20:59\,$\to$\,\textbf{20:58} & -- & $-$1 & +1.0 & na EIP 3503 (21:03) \\
1 & 3 & Księżnej Dobrawy & -- & 21:06\,$\to$\,\textbf{21:07} & +1 & +1.0 & z R 90113 (21:03) \\
2 & 3 & Księżnej Dobrawy & -- & 5:22\,$\to$\,\textbf{5:27} & +5 & +0.0 & \textit{zmiana towarzysząca} \\
2 & 2 & Iława ul. Lipowy Dwór (ogrody) & -- & 5:42\,$\to$\,\textbf{5:40} & $-$2 & +0.0 & \textit{zmiana towarzysząca} \\
2 & 3 & Nowa Wieś & -- & 6:02\,$\to$\,\textbf{5:58} & $-$4 & +1.0 & z R 90205 (05:44) \\
2 & 2 & Dworcowa & 6:20\,$\to$\,\textbf{6:17} & -- & $-$3 & +1.0 & na R 90745 (06:22) \\
2 & 2 & Iława ul. Lipowy Dwór (ogrody) & -- & 7:11\,$\to$\,\textbf{7:09} & $-$2 & +1.0 & z EIP 5300 (06:52) \\
2 & 2 & Iława ul. Lipowy Dwór (ogrody) & -- & 7:40\,$\to$\,\textbf{7:36} & $-$4 & +1.0 & z R 90240 (07:22) \\
2 & 3 & Księżnej Dobrawy & -- & 8:26\,$\to$\,\textbf{8:25} & $-$1 & +0.0 & \textit{zmiana towarzysząca} \\
2 & 3 & Dworcowa & 9:15\,$\to$\,\textbf{9:10} & -- & $-$5 & +2.0 & na R 90751 (09:15); na IC 7503 (09:17) \\
2 & 3 & Księżnej Dobrawy & -- & 9:45\,$\to$\,\textbf{9:42} & $-$3 & +0.0 & \textit{zmiana towarzysząca} \\
2 & 3 & Dworcowa & 10:30\,$\to$\,\textbf{10:25} & -- & $-$5 & +1.0 & na IC 5710 (10:30) \\
2 & 3 & Księżnej Dobrawy & -- & 10:35\,$\to$\,\textbf{10:33} & $-$2 & +0.0 & \textit{zmiana towarzysząca} \\
2 & 3 & Dworcowa & 11:22\,$\to$\,\textbf{11:19} & -- & $-$3 & +1.0 & na IC 1547 (11:24) \\
2 & 3 & Dworcowa & 12:22\,$\to$\,\textbf{12:26} & -- & +4 & +1.0 & na IC 1706 (12:40) \\
2 & 3 & Dworcowa & 12:41\,$\to$\,\textbf{12:44} & -- & +3 & +1.0 & na IC 5114 (12:59) \\
2 & 3 & Nowa Wieś & -- & 13:35\,$\to$\,\textbf{13:33} & $-$2 & +1.0 & z IC 7515 (13:16) \\
2 & 2 & Dworcowa & 13:51\,$\to$\,\textbf{13:48} & -- & $-$3 & +1.0 & na EIP 5404 (13:53) \\
2 & 3 & Dworcowa & 14:57\,$\to$\,\textbf{14:54} & -- & $-$3 & +1.0 & na IC 8114 (14:59) \\
2 & 3 & Nowa Wieś & -- & 17:21\,$\to$\,\textbf{17:18} & $-$3 & +1.0 & z EIC 4505 (17:02) \\
2 & 3 & Dworcowa & 18:07\,$\to$\,\textbf{18:02} & -- & $-$5 & +0.0 & \textit{zmiana towarzysząca} \\
2 & 2 & Dworcowa & 18:51\,$\to$\,\textbf{18:53} & -- & +2 & +1.0 & na IC 1517 (19:08) \\
3 & 2 & Iława ul. Lipowy Dwór (ogrody) & -- & 15:10\,$\to$\,\textbf{15:08} & $-$2 & +0.0 & z EIP 5304 (14:52); \textit{traci: z EIP 3553 (15:02)} \\
3 & 2 & Dworcowa & 15:52\,$\to$\,\textbf{15:48} & -- & $-$4 & +1.0 & na EIP 5306 (15:53) \\
3 & 1 & Dworcowa & 16:27\,$\to$\,\textbf{16:25} & -- & $-$2 & +0.0 & \textit{zmiana towarzysząca} \\
\end{longtable}

\subsection*{Niedziela i święta}
Plan na dzień 4.10.2026, sprawdzony na 4 dniach tego typu (4.10.2026--25.10.2026). Dobre przesiadki dziennie (pary pociąg--linia, z zapasem na typowe opóźnienie pociągu): obecnie 46.

\begin{tabular}{@{}lrrr@{}}\toprule
Etap & Zmienione kursy (razem) & Dobre przesiadki dziennie & Zmiana \\ \midrule
1 & 20 & 70 & +52\,\% \\
2 & 28 & 76 & +64\,\% \\
\bottomrule\end{tabular}

\begin{longtable}{@{}c c >{\raggedright}p{2.4cm} l l r r >{\raggedright\arraybackslash}p{4.3cm}@{}}\toprule
Etap & Linia & Kierunek & Przyjazd & Odjazd & Inne & Zysk & Pociągi \\
 & & & \multicolumn{2}{c}{na Dworcowej} & przyst. & /dzień & (dobra przesiadka po zmianie) \\ \midrule\endhead
\bottomrule\endfoot
1 & 2 & Dworcowa & 7:45\,$\to$\,\textbf{7:48} & -- & +3 & +1.0 & na EIP 1501 (08:03) \\
1 & 3 & Dworcowa & 8:44\,$\to$\,\textbf{8:48} & -- & +4 & +2.0 & na IC 5116 (09:00); na EIP 1503 (09:03) \\
1 & 3 & Nowa Wieś & -- & 9:16\,$\to$\,\textbf{9:14} & $-$2 & +2.0 & z IC 1553 (08:57); z IC 5116 (08:58) \\
1 & 1 & Długa & -- & 9:35\,$\to$\,\textbf{9:32} & $-$3 & +1.0 & z IC 7503 (09:16) \\
1 & 3 & Dworcowa & 10:46\,$\to$\,\textbf{10:48} & -- & +2 & +1.0 & na EIP 4509 (11:03) \\
1 & 1 & Dworcowa & 10:47\,$\to$\,\textbf{10:48} & -- & +1 & +1.0 & na EIP 4509 (11:03) \\
1 & 1 & Długa & -- & 10:58\,$\to$\,\textbf{11:03} & +5 & +2.0 & z EIP 5402 (10:52); z IC 1513 (10:55) \\
1 & 3 & Księżnej Dobrawy & -- & 11:00\,$\to$\,\textbf{11:03} & +3 & +1.0 & z IC 1513 (10:55) \\
1 & 3 & Księżnej Dobrawy & -- & 12:20\,$\to$\,\textbf{12:19} & $-$1 & +1.0 & z EIP 3509 (12:02) \\
1 & 3 & Dworcowa & 12:37\,$\to$\,\textbf{12:38} & -- & +1 & +1.0 & na EIP 5352 (12:53) \\
1 & 2 & Iława ul. Lipowy Dwór (ogrody) & -- & 14:38\,$\to$\,\textbf{14:42} & +4 & +2.0 & z IC 5714 (14:33); z R 55205 (14:38) \\
1 & 3 & Dworcowa & 15:22\,$\to$\,\textbf{15:24} & -- & +2 & +1.0 & na R 90244 (15:39) \\
1 & 3 & Nowa Wieś & -- & 15:35\,$\to$\,\textbf{15:32} & $-$3 & +1.0 & z IC 7107 (15:16) \\
1 & 1 & Długa & -- & 15:40\,$\to$\,\textbf{15:42} & +2 & +1.0 & z R 90244 (15:38) \\
1 & 1 & Dworcowa & 16:57\,$\to$\,\textbf{16:58} & -- & +1 & +1.0 & na R 90636 (17:13) \\
1 & 3 & Księżnej Dobrawy & -- & 17:21\,$\to$\,\textbf{17:19} & $-$2 & +1.0 & z EIC 4505 (17:02) \\
1 & 2 & Iława ul. Lipowy Dwór (ogrody) & -- & 18:10\,$\to$\,\textbf{18:09} & $-$1 & +1.0 & z EIP 5406 (17:52) \\
1 & 2 & Dworcowa & 18:51\,$\to$\,\textbf{18:53} & -- & +2 & +1.0 & na IC 1517 (19:08) \\
1 & 3 & Nowa Wieś & -- & 19:15\,$\to$\,\textbf{19:14} & $-$1 & +1.0 & z IC 1517 (18:57) \\
1 & 3 & Dworcowa & 20:02\,$\to$\,\textbf{20:00} & -- & $-$2 & +1.0 & na IC 5320 (20:05) \\
2 & 3 & Nowa Wieś & -- & 6:05\,$\to$\,\textbf{6:04} & $-$1 & +0.0 & \textit{zmiana towarzysząca} \\
2 & 3 & Dworcowa & 6:51\,$\to$\,\textbf{6:48} & -- & $-$3 & +1.0 & na EIP 5300 (06:53) \\
2 & 3 & Księżnej Dobrawy & -- & 7:58\,$\to$\,\textbf{8:02} & +4 & +1.0 & z IC 1521 (07:55) \\
2 & 1 & Długa & -- & 8:15\,$\to$\,\textbf{8:12} & $-$3 & +1.0 & z IC 1521 (07:55) \\
2 & 3 & Księżnej Dobrawy & -- & 10:00\,$\to$\,\textbf{10:03} & +3 & +0.0 & \textit{zmiana towarzysząca} \\
2 & 2 & Iława ul. Lipowy Dwór (ogrody) & -- & 10:51\,$\to$\,\textbf{10:46} & $-$5 & +1.0 & z IC 5710 (10:29) \\
2 & 3 & Dworcowa & 11:52\,$\to$\,\textbf{11:48} & -- & $-$4 & +1.0 & na EIP 5302 (11:53) \\
2 & 2 & Dworcowa & 15:18\,$\to$\,\textbf{15:16} & -- & $-$2 & +0.8 & na R 90634 (15:21) \\
\end{longtable}

}
\IfFileExists{\DrukPrzesiadki}{\includepdf[pages=-]{\DrukPrzesiadki}}{\typeout{Brak pliku \DrukPrzesiadki: uruchom print_reports.py --pdf}}
\IfFileExists{\DrukKorekty}{\includepdf[pages=-]{\DrukKorekty}}{\typeout{Brak pliku \DrukKorekty: uruchom print_reports.py --pdf}}}
